\notebook{06-simulating-hawkes}
The questions of this section concern the simulation of Section~\ref{notes:sec.simulation}:
thinning, the update of the state,
and the exact scheme of Theorem~\ref{notes:thm.dassiosZhao},
with the simulators of the package.
A process of two types is read as the market orders,
type 0 the market buys and type 1 the market sells,
since the package counts types from 0.

\begin{question}
\label{question.thinningBoundAfterTheJump}
Algorithm~\ref{notes:algo.ogata} draws candidates at the rate of a bound,
and keeps each with probability the total intensity $\totalIntensity$ over the bound.
The two functions below run it on the state $\decayedCounts$ of a pair of market orders,
and return the number of events per second up to \codehl{horizon}.
\codehl{rng.exponential(scale)} draws an exponential of mean \codehl{scale},
and \codehl{rng.choice(2, p=p)} draws 0 or 1 with probabilities \codehl{p}.
\begin{lstlisting}[caption={Alternative A}, label=listing.thinningBoundAfterTheJumpA]
def thinning_rate(params, horizon, rng):
    mu_bar = params.baseline.sum()
    jumps = params.excitation.sum(axis=0)
    time, state, count = 0.0, np.zeros(2), 0
    bound = mu_bar
    while True:
        wait = rng.exponential(1 / bound)
        time += wait
        if time > horizon:
            return count / horizon
        state *= np.exp(-params.decay * wait)
        intensity = mu_bar + jumps @ state
        if rng.random() * bound <= intensity:
            lam = params.baseline + params.excitation @ state
            state[rng.choice(2, p=lam / lam.sum())] += 1.0
            count += 1
            bound = mu_bar + jumps @ state
\end{lstlisting}
Alternative B is Alternative A with the last five lines of its loop replaced by the following.
\begin{lstlisting}[caption={Alternative B, the end of its loop}, label=listing.thinningBoundAfterTheJumpB]
        if rng.random() * bound <= intensity:
            lam = params.baseline + params.excitation @ state
            state[rng.choice(2, p=lam / lam.sum())] += 1.0
            count += 1
        bound = intensity
\end{lstlisting}
Each is run on the self-exciting pair,
\begin{equation*}
 \baseIntensity = (0.25, 0.25),
 \qquad
 \excitation = \begin{pmatrix} 2.8 & 0.2 \\ 0.2 & 2.8 \end{pmatrix},
 \qquad
 \decay = 4,
\end{equation*}
over $2 \times 10^{4}$ seconds,
where the standard error of the rate is $0.04$.
\begin{lstlisting}[caption={Run once with each alternative}, label=listing.thinningBoundAfterTheJumpRun]
params = config.self_exciting_pair_params()
rng = np.random.default_rng(8)
print(round(thinning_rate(params, 2e4, rng), 2))
\end{lstlisting}
Which statement holds?
\begin{choices}
 \choice Both print a number close to 2.
 \choice A prints a number clearly below 2, and B a number close to 2.
 \choice A prints a number close to 2, and B a number clearly above 2.
 \correctchoice A prints a number close to 2, and B a number clearly below 2.
\end{choices}
\explanation{The columns of $\excitation/\decay$ sum to $0.75$, so $\totalRate = 0.5/0.25 = 2$.
A bounds by the post-jump value $\totalIntensity(\arrivalTimes[]_{n}+)$ until the next event,
which dominates $\totalIntensity$ since $\excitation \geq 0$:
it is exact, and wastes only its rejections.
B bounds by the pre-jump value, smaller by the jump of 3,
so its candidates come at a rate below $\totalIntensity$ and part of the excitation is lost.}
\end{question}

\begin{question}
\label{question.recordedStatesAliasTheJump}
The update~\eqref{notes:eq.update_of_exponential_intensity} carries the state $\decayedCounts$
from one market order to the next,
and Remark~\ref{notes:remark.prePostJump} tells its two values apart:
the pre-jump value $\decayedCounts(\arrivalTimes[]_{n})$ enters the intensity,
and the post-jump value $\decayedCounts(\arrivalTimes[]_{n}+)$ is carried forward.
Each function below is written to record $\decayedCounts(\arrivalTimes[]_{n})$ at every event,
and the two differ in one line.
\begin{lstlisting}[caption={Alternative A}, label=listing.recordedStatesAliasTheJumpA]
def recorded_states(decay, times, types):
    state, previous, record = np.zeros(2), 0.0, []
    for now, event_type in zip(times, types):
        state = state * np.exp(-decay * (now - previous))
        record.append(state)
        state[event_type] += 1.0
        previous = now
    return record
\end{lstlisting}
\begin{lstlisting}[caption={Alternative B}, label=listing.recordedStatesAliasTheJumpB]
def recorded_states(decay, times, types):
    state, previous, record = np.zeros(2), 0.0, []
    for now, event_type in zip(times, types):
        state *= np.exp(-decay * (now - previous))
        record.append(state)
        state[event_type] += 1.0
        previous = now
    return record
\end{lstlisting}
A market buy, a market sell and a market buy arrive half a second apart,
and at the decay $\decay = 2\log 2$ the state halves in half a second.
\begin{lstlisting}[caption={Run once with each alternative}, label=listing.recordedStatesAliasTheJumpRun]
times, types = [0.5, 1.0, 1.5], [0, 1, 0]
record = recorded_states(2 * np.log(2), times, types)
print([state.tolist() for state in record])
\end{lstlisting}
What do the two alternatives print?
\begin{choices}
 \choice A prints \codehl{[[0.0, 0.0], [0.5, 0.0], [0.25, 0.5]]},\\
  and B prints \codehl{[[1.25, 0.5], [1.25, 0.5], [1.25, 0.5]]}.
 \choice Both print \codehl{[[0.0, 0.0], [0.5, 0.0], [0.25, 0.5]]}.
 \correctchoice A prints \codehl{[[1.0, 0.0], [0.5, 1.0], [1.25, 0.5]]},\\
  and B prints \codehl{[[1.25, 0.5], [1.25, 0.5], [1.25, 0.5]]}.
 \choice Both print \codehl{[[1.0, 0.0], [0.5, 1.0], [1.25, 0.5]]}.
\end{choices}
\explanation{In A, \codehl{state * ...} builds a new array, which the record holds
and to which the jump then adds:
every entry is the post-jump $\decayedCounts(\arrivalTimes[]_{n}+)$.
In B, \codehl{*=} decays one array in place, and the record holds it three times, at its final value.
Recording \codehl{state.copy()} gives the pre-jump states,
\codehl{[[0.0, 0.0], [0.5, 0.0], [0.25, 0.5]]}.}
\end{question}

\begin{question}
\label{question.intensityByRecursionAgainstSum}
The total intensity $\totalIntensity(\arrivalTimes[]_{n})$ standing at each market order of a path
is computed in two ways.
Alternative A sums over the events before $\arrivalTimes[]_{n}$, as in~\eqref{notes:eq.decayedCounts};
Alternative B carries the state by the update~\eqref{notes:eq.update_of_exponential_intensity}.
\codehl{times} and \codehl{types} are arrays of the $n$ times and types of the path.
\begin{lstlisting}[caption={Alternative A}, label=listing.intensityByRecursionAgainstSumA]
def total_intensities_a(params, times, types):
    jumps = params.excitation.sum(axis=0)
    out = np.empty(times.size)
    for n, now in enumerate(times):
        weights = np.exp(-params.decay * (now - times[:n]))
        out[n] = params.baseline.sum() + np.sum(weights * jumps[types[:n]])
    return out
\end{lstlisting}
\begin{lstlisting}[caption={Alternative B}, label=listing.intensityByRecursionAgainstSumB]
def total_intensities_b(params, times, types):
    jumps = params.excitation.sum(axis=0)
    state, previous = np.zeros(2), 0.0
    out = np.empty(times.size)
    for n, (now, event_type) in enumerate(zip(times, types)):
        state *= np.exp(-params.decay * (now - previous))
        out[n] = params.baseline.sum() + jumps @ state
        state[event_type] += 1.0
        previous = now
    return out
\end{lstlisting}
Both are timed on the first $n$ events of one path,
for $n$ from 1000 to 16000, doubling.
Which statement holds?
\begin{choices}
 \correctchoice The time per event of B stays about the same and that of A grows with $n$:
  B is faster at every $n$, by a factor that grows with $n$.
 \choice The time per event of both stays about the same,
  and B is faster by a constant factor.
 \choice A is faster at every $n$, its loop body being one vectorised expression.
 \choice The time per event of both grows with $n$,
  since each loops over the events of the path.
\end{choices}
\explanation{The body of A costs a term for every earlier event,
so the path costs $O(n^{2})$,
and vectorising the sum lowers the constant and not the growth.
The update costs $O(1)$ per event.}
\end{question}

\begin{question}
\label{question.exactWaitOnGivenUniforms}
The step of the exact scheme of the package, Algorithm~\ref{notes:algo.dassiosZhao}, is quoted below.
In it, \codehl{self.total\_intensity} is $\totalIntensity$ at the current state,
and \codehl{self.\_baseline\_total} is $\bar{\baseIntensity}$.
The call \codehl{self.\_advance(w)} moves the clock by \codehl{w}
and multiplies the state by $e^{-\decay w}$,
and the call \codehl{self.\_draw\_type\_and\_jump()} draws the type
and adds one to its coordinate of the state.
\begin{lstlisting}[caption={\protect\codehl{unito26.lob.hawkes.ExponentialHawkes.step}, its comments elided.}, label=listing.exactWaitStep]
def step(self) -> tuple[float, int]:
    beta = self.params.decay
    excess = self.total_intensity - self._baseline_total  # D >= 0

    u1, u2 = self.rng.random(2)

# ...
    wait_baseline = -np.log(u1) / self._baseline_total

# ...
    if excess <= 0.0:
        wait_excited = np.inf
    else:
        remaining = 1.0 + beta * np.log(u2) / excess
        wait_excited = -np.log(remaining) / beta if remaining > 0.0 else np.inf

    self._advance(min(wait_baseline, wait_excited))
    return self.time, self._draw_type_and_jump()
\end{lstlisting}
The simulator is started just after a market order of the asymmetric pair,
\begin{equation*}
 \baseIntensity = (0.2, 0.6),
 \qquad
 \excitation = \begin{pmatrix} 1.2 & 4.8 \\ 0.4 & 0.8 \end{pmatrix},
 \qquad
 \decay = 4,
\end{equation*}
from the state $\decayedCounts(0+) = (1.5, 1.0)$.
\codehl{Uniforms(values)} stands in for the generator:
its \codehl{random(2)} returns the next two of \codehl{values} as an array,
and its \codehl{random()} the next one.
Take $\log 2 \approx 0.693$.
\begin{lstlisting}[caption={One step on given uniforms, in two cases}, label=listing.exactWaitRun]
params = config.asymmetric_pair_params()
for u2 in (np.exp(-1.0), np.exp(-3.0)):
    simulator = ExponentialHawkes.from_state(params, [1.5, 1.0])
    simulator.rng = Uniforms([np.exp(-0.4), u2, 0.75])
    time, event_type = simulator.step()
    print(round(time, 3))
\end{lstlisting}
What does the snippet print?
\begin{choices}
 \choice \codehl{0.5}, then \codehl{0.5}.
 \choice \codehl{0.5}, then \codehl{inf}.
 \choice \codehl{0.173}, then \codehl{0.32}.
 \correctchoice \codehl{0.173}, then \codehl{0.5}.
 \choice \codehl{0.173}, then \codehl{inf}.
\end{choices}
\explanation{By~\eqref{notes:eq.dzDraws}, with $D = 1.6 \cdot 1.5 + 5.6 \cdot 1.0 = 8$,
$\tau_1 = 0.4/0.8 = 0.5$,
and $\tau_2 = -\tfrac{1}{4}\log(1 - \tfrac{4}{8}) = \tfrac{1}{4}\log 2 \approx 0.173$ in the first case;
in the second $1 - \tfrac{12}{8} < 0$, so $\tau_2 = +\infty$ and the baseline clock fires at $0.5$.
The column sums of $\branchingMatrix$ taken for those of $\excitation$ give $D = 2$
and an infinite $\tau_2$ in both cases.}
\end{question}

\begin{question}
\label{question.typeDrawnBeforeTheDecay}
Equation~\eqref{notes:eq.dzType} draws the type of the next market order
on the intensities at its arrival, before its jump.
The package inverts their cumulative sum, as below;
\codehl{np.searchsorted(cumulative, threshold)} is the first index
at which \codehl{cumulative} reaches \codehl{threshold}.
\begin{lstlisting}[caption={\protect\codehl{unito26.lob.hawkes.\_HawkesState}: the intensities, the decay and the type draw, their docstrings elided.}, label=listing.typeDrawState]
@property
def intensities(self) -> np.ndarray:
# ...
    return self.params.baseline + self.params.excitation @ self.decayed_counts
# ...
def _advance(self, elapsed: float) -> None:
# ...
    self.time += elapsed
    self.decayed_counts *= np.exp(-self.params.decay * elapsed)

def _draw_type_and_jump(self) -> int:
# ...
    intensities = self.intensities
    cumulative = np.cumsum(intensities)
    threshold = self.rng.random() * cumulative[-1]
# ...
    event_type = min(int(np.searchsorted(cumulative, threshold)), self.params.dimension - 1)
    self.decayed_counts[event_type] += 1.0
    return event_type
\end{lstlisting}
The class \codehl{DrawFirst} is \codehl{ExponentialHawkes}
with the last two lines of its step, Listing~\ref{listing.exactWaitStep},
replaced by the three below.
\begin{lstlisting}[caption={The end of the step of \protect\codehl{DrawFirst}}, label=listing.typeDrawnBeforeTheDecayVariant]
event_type = self._draw_type_and_jump()
self._advance(min(wait_baseline, wait_excited))
return self.time, event_type
\end{lstlisting}
Both start from the state of Question~\ref{question.exactWaitOnGivenUniforms},
on the uniforms of its first case.
\begin{lstlisting}[caption={One step of each class}, label=listing.typeDrawnBeforeTheDecayRun]
params = config.asymmetric_pair_params()
for simulator_class in (ExponentialHawkes, DrawFirst):
    simulator = simulator_class.from_state(params, [1.5, 1.0])
    simulator.rng = Uniforms([np.exp(-0.4), np.exp(-1.0), 0.75])
    time, event_type = simulator.step()
    print(event_type, simulator.decayed_counts.round(3).tolist())
\end{lstlisting}
What does the snippet print?
\begin{choices}
 \choice \codehl{1 [0.75, 1.5]} on both lines.
 \choice \codehl{1 [0.75, 1.5]}, then \codehl{0 [1.75, 0.5]}.
 \choice \codehl{1 [0.75, 1.5]}, then \codehl{0 [2.5, 1.0]}.
 \choice \codehl{0 [1.75, 0.5]}, then \codehl{0 [1.25, 0.5]}.
 \correctchoice \codehl{1 [0.75, 1.5]}, then \codehl{0 [1.25, 0.5]}.
\end{choices}
\explanation{The wait multiplies the state by $1 - \tfrac{4}{8} = \tfrac{1}{2}$:
at the arrival the intensities are $(3.5, 1.3)$,
and the threshold $0.75 \cdot 4.8 = 3.6$ lies past the buys, so the order is a market sell.
\codehl{DrawFirst} draws on the intensities $(6.8, 2.0)$ left by the last event,
where $6.6$ falls among the buys,
and then halves the state with its own jump in it, from $(2.5, 1.0)$ to $(1.25, 0.5)$.}
\end{question}

\begin{question}
\label{question.clippedExcitationClockAddsEvents}
The excitation clock of the exact scheme is defective, Remark~\ref{notes:remark.defectiveClock}:
with probability $e^{-D/\decay}$ it never fires.
The class \codehl{Clipped} is \codehl{ExponentialHawkes}
with the two lines under \codehl{else} in its step, Listing~\ref{listing.exactWaitStep},
replaced by the two below, which never return infinity.
\begin{lstlisting}[caption={The excitation clock of \protect\codehl{Clipped}}, label=listing.clippedExcitationClockAddsEventsVariant]
remaining = max(1.0 + beta * np.log(u2) / excess, 1e-12)
wait_excited = -np.log(remaining) / beta
\end{lstlisting}
The flow below decays at 50 per second,
so that its excitation halves in about 14 milliseconds,
a time scale of the market orders of a liquid stock.
Take $\log 10^{-12} \approx -27.6$.
The snippet is run with \codehl{simulator\_class} each of the two classes,
over $2 \times 10^{4}$ seconds, where the standard error of the rate is $0.02$.
\begin{lstlisting}[caption={Run once with each class}, label=listing.clippedExcitationClockAddsEventsRun]
params = HawkesParams(
    baseline=[0.5, 0.5],
    excitation=[[15.0, 10.0], [10.0, 15.0]],
    decay=50.0,
)
events = list(simulator_class(params, rng=4).events(2e4))
print(round(len(events) / 2e4, 2))
\end{lstlisting}
Which statement holds?
\begin{choices}
 \choice Both print a number close to 2.
 \correctchoice \codehl{ExponentialHawkes} prints a number close to 2,
  and \codehl{Clipped} a number clearly above 2.
 \choice \codehl{ExponentialHawkes} prints a number close to 2,
  and \codehl{Clipped} a number clearly below 2.
 \choice \codehl{ExponentialHawkes} prints a number clearly below 2,
  its excitation clock failing to fire on some waits,
  and \codehl{Clipped} a number close to 2.
\end{choices}
\explanation{The columns of $\excitation/\decay$ sum to $0.5$, so $\totalRate = 1/0.5 = 2$.
\codehl{Clipped} makes every defective clock fire at $27.6/50 \approx 0.55$ seconds,
before the baseline clock, of mean $1/\bar{\baseIntensity} = 1$ second, more often than not:
each such firing is an event the model does not have, and it excites further events.}
\end{question}

\begin{question}
\label{question.selfAndCrossExcitingShareTheirTimes}
The self-exciting and the cross-exciting pairs of the package share
$\baseIntensity = (0.25, 0.25)$ and $\decay = 4$;
their excitation matrices are
\begin{equation*}
 \begin{pmatrix} 2.8 & 0.2 \\ 0.2 & 2.8 \end{pmatrix}
 \quad\text{and}\quad
 \begin{pmatrix} 0.2 & 2.8 \\ 2.8 & 0.2 \end{pmatrix}.
\end{equation*}
In the first a market buy mostly excites market buys,
in the second market sells.
Both are simulated by the step of Listing~\ref{listing.exactWaitStep},
which reads the total intensity as below,
and whose type draw takes one uniform.
\begin{lstlisting}[caption={\protect\codehl{unito26.lob.hawkes.\_HawkesState}: the total intensity.}, label=listing.selfAndCrossTotalIntensity]
# ...
    self._total_jump = params.excitation.sum(axis=0)
# ...
@property
def total_intensity(self) -> float:
# ...
    return self._baseline_total + float(self._total_jump @ self.decayed_counts)
\end{lstlisting}
Each path below is an array with a row (time, type) per event.
\codehl{np.allclose(x, y)} raises \codehl{ValueError} when \codehl{x} and \codehl{y}
differ in length, neither being of length one.
\begin{lstlisting}[caption={The two pairs on one seed}, label=listing.selfAndCrossExcitingShareTheirTimesRun]
self_exciting = config.self_exciting_pair_params()
cross_exciting = config.cross_exciting_pair_params()
paths = [
    np.array(list(ExponentialHawkes(params, rng=3).events(100.0)))
    for params in (self_exciting, cross_exciting)
]
times = [path[:, 0] for path in paths]
types = [path[:, 1] for path in paths]
print(np.allclose(times[0], times[1]), np.array_equal(types[0], types[1]))
\end{lstlisting}
What does the snippet print?
\begin{choices}
 \choice \codehl{True True}
 \choice \codehl{False False}
 \correctchoice \codehl{True False}
 \choice \codehl{ValueError}
\end{choices}
\explanation{Both matrices have column sums $(3, 3)$,
so $\totalIntensity = 0.5 + 3\,\mathbf{1}^{\transpose}\decayedCounts$
depends on the times of the past events and not on their types:
the same uniforms give the same waits, and the times agree to rounding.
The type is drawn on the full intensities, where the two matrices differ,
Remark~\ref{notes:remark.typeDraw}.}
\end{question}

\begin{question}
\label{question.eventsHoldTheCrossingEvent}
A simulation is often run in two stretches,
a warm-up to a time and then the run proper.
The simulators yield their market orders up to a horizon by the generator quoted below;
\codehl{self.\_crossing} is \codehl{None} when a simulator is built,
and \codehl{self.step()} draws the next event, advancing the clock and the state.
\begin{lstlisting}[caption={\protect\codehl{unito26.lob.hawkes.\_HawkesState.events}, its docstring elided.}, label=listing.eventsHoldTheCrossingEventQuote]
def events(self, horizon: float) -> Iterator[tuple[float, int]]:
# ...
    while True:
        if self._crossing is not None:
            time, event_type = self._crossing
            self._crossing = None
        else:
            time, event_type = self.step()
        if time > horizon:
            self._crossing = (time, event_type)
            return
        yield time, event_type
\end{lstlisting}
The class \codehl{Dropping} overrides it.
The one-way pair has a total baseline $\bar{\baseIntensity} = 0.5$,
so that thirty seconds pass without an event with a probability below $e^{-15}$.
\begin{lstlisting}[caption={A simulator that drops the event past the horizon}, label=listing.eventsHoldTheCrossingEventVariant]
class Dropping(ExponentialHawkes):
    def events(self, horizon):
        while True:
            time, event_type = self.step()
            if time > horizon:
                return
            yield time, event_type
\end{lstlisting}
\begin{lstlisting}[caption={Two stretches against one run}, label=listing.eventsHoldTheCrossingEventRun]
params = config.one_way_pair_params()
for simulator_class in (ExponentialHawkes, Dropping):
    simulator = simulator_class(params, rng=11)
    first = list(simulator.events(10.0))
    rest = list(simulator.events(40.0))
    whole = list(simulator_class(params, rng=11).events(40.0))
    print(first + rest == whole, len(whole) - len(first + rest))
\end{lstlisting}
What does the snippet print?
\begin{choices}
 \correctchoice \codehl{True 0}, then \codehl{False 1}.
 \choice \codehl{True 0} on both lines.
 \choice \codehl{False 1} on both lines.
 \choice \codehl{False} and a negative number, on both lines.
\end{choices}
\explanation{The first event past 10 seconds is drawn, and its draw has already moved the clock and the state.
The package holds it and yields it first in the second call,
so that the two stretches are the single run;
\codehl{Dropping} discards it, and loses that event alone.
A second call continues from where the first stopped, not from time zero.}
\end{question}

\begin{question}
\label{question.seedsOfAnOrderFlowMonteCarlo}
A Monte Carlo average over paths of the order flow needs paths that are independent,
and a study that can be rerun needs them reproducible.
The simulators take an integer seed or a generator, as below.
\begin{lstlisting}[caption={\protect\codehl{unito26.lob.hawkes.\_HawkesState.\_\_init\_\_}: the generator.}, label=listing.seedsOfAnOrderFlowMonteCarloQuote]
# ...
self.rng = rng if isinstance(rng, np.random.Generator) else np.random.default_rng(rng)
# ...
\end{lstlisting}
Two generators built by \codehl{np.random.default\_rng(5)} draw the same numbers.
The weakly exciting pair has a total rate of 2 events a second,
so that a path of ten seconds is empty with a negligible probability.
\begin{lstlisting}[caption={Paths from a seed and from a generator}, label=listing.seedsOfAnOrderFlowMonteCarloRun]
params = config.weakly_exciting_pair_params()
rng = np.random.default_rng(5)

def path(seed_or_generator):
    simulator = ExponentialHawkes(params, rng=seed_or_generator)
    return list(simulator.events(10.0))

print(path(5) == path(5), path(rng) == path(rng), path(5) == path(rng))
\end{lstlisting}
What does the snippet print?
\begin{choices}
 \choice \codehl{True False True}
 \correctchoice \codehl{True False False}
 \choice \codehl{True True True}
 \choice \codehl{False False False}
\end{choices}
\explanation{An integer seed builds a fresh generator, so \codehl{path(5)} is one path at every call.
A generator is used as given and advanced by every draw:
the calls of \codehl{path(rng)} read consecutive stretches of its stream,
and by the third it has moved past the numbers with which \codehl{default\_rng(5)} starts.
Paths that must be independent and reproducible are drawn from one generator,
seeded once and passed to every simulator.}
\end{question}

\begin{question}
\label{question.transposedKernelChangesTheRate}
The records \codehl{config.ASYMMETRIC\_PAIR\_PARAMS}
list the entries of the excitation matrix of the asymmetric pair in the order below,
\codehl{Component} being the type excited, $e$,
and \codehl{Cause} the type exciting it, $e'$,
as in Remark~\ref{notes:remark.indexOrder}.
\begin{center}
\begin{tabular}{ccc}
 \toprule
 \codehl{Component} & \codehl{Cause} & \codehl{Kernel} \\
 \midrule
 0 & 0 & 1.2 \\
 0 & 1 & 4.8 \\
 1 & 0 & 0.4 \\
 1 & 1 & 0.8 \\
 \bottomrule
\end{tabular}
\end{center}
The baselines are $\baseIntensity = (0.2, 0.6)$ and the decay is $\decay = 4$.
The snippet reads the matrix off the records in two ways,
and simulates each specification over $2 \times 10^{4}$ seconds,
where the standard errors of the two rates are below $0.04$.
\begin{lstlisting}[caption={The records read row by row, and transposed}, label=listing.transposedKernelChangesTheRateRun]
records = config.ASYMMETRIC_PAIR_PARAMS
kernel = np.array([record["Kernel"] for record in records])
for excitation in (kernel.reshape(2, 2), kernel.reshape(2, 2).T):
    params = HawkesParams(
        baseline=[0.2, 0.6], excitation=excitation, decay=4.0
    )
    events = list(ExponentialHawkes(params, rng=1).events(2e4))
    print(round(len(events) / 2e4, 2))
\end{lstlisting}
What does the snippet print?
\begin{choices}
 \choice A number close to 3, twice.
 \correctchoice A number close to 3, then a number close to 2.
 \choice A number close to 2, then a number close to 3.
 \choice A number close to 2, twice.
\end{choices}
\explanation{\codehl{reshape} fills rows first, so the first matrix is $\excitation$, row excited.
With $\branchingMatrix = \excitation/4$, $\det(I - \branchingMatrix) = 0.44$,
and $\stationaryIntensity = (I - \branchingMatrix)^{-1}\baseIntensity$
of~\eqref{notes:eq.stationaryIntensity} is $(2, 1)$; with the transpose it is $(0.5, 1.5)$.
Both have $\branchingRatio = 0.6$ and pass the condition of Theorem~\ref{notes:thm.stationarity}:
at a fixed baseline the transpose changes the rate.
The rate $\bar{\baseIntensity}/(1 - \branchingRatio) = 2$ holds for the transpose,
of which $\baseIntensity$ is a right Perron vector,
and not for $\excitation$.}
\end{question}
